\documentclass[11pt,a4paper]{article}
\usepackage[T2A]{fontenc}
\usepackage[utf8]{inputenc}
\usepackage[russian,english]{babel}
\usepackage{amsmath,amssymb,amsthm,mathtools}
\usepackage{setspace}
\usepackage{enumitem}
\usepackage[hidelinks]{hyperref}
\setstretch{1.3}
\frenchspacing
\hyphenpenalty=10000

\newcommand{\head}[1]{\par\noindent\textbf{\underline{#1}}\quad}
\newcommand{\sheettitle}[4]{% title, subtitle, homework-flag, date
  \begin{center}
  {\huge\bfseries #1\par}\vspace{\dimexpr1.2em-1pt\relax minus 3.6pt}
  {\Large #2\par}\vspace{\dimexpr1em+1pt\relax minus 1.2pt}
  \if\relax\detokenize{#3}\relax\else{\Large #3\par}\vspace{\dimexpr1.15em-1pt\relax minus 3.6pt}\fi
  {\foreignlanguage{russian}{#4}\par}
  \end{center}\vspace{0.5pt minus 1.2pt}}
\newcommand{\src}[1]{\unskip\nobreak\hfil\penalty50\hskip1em\hbox{}\nobreak\hfill(#1){\parfillskip=0pt\par}}
\DeclareMathOperator{\tr}{tr}
\DeclareMathOperator{\adj}{adj}
\newcommand{\T}{\mathsf T}
\DeclareMathOperator{\rank}{rank}
\begin{document}
\sheettitle{Linear algebra -- 10}{Matrix norms and localisation of the spectrum}{Homework}{12 сентября 2026 года}

\head{Theoretical minimum.} Operator and Frobenius norms; the spectral radius and Gelfand's
formula; Gershgorin discs; Courant--Fischer and Cauchy interlacing; Weyl's inequalities
$\lambda_{i+j-1}(A+B)\leqslant\lambda_i(A)+\lambda_j(B)$; the Hoffman--Wielandt inequality; the
numerical range $W(A)=\{\langle Ax,x\rangle:\|x\|=1\}$ and the condition number
$\varkappa(A)=\|A\|\,\|A^{-1}\|$.

\par\vspace{-1pt}\noindent\rule{\linewidth}{0.4pt}\par\vspace{7pt}

\begin{enumerate}
\item Prove Gershgorin's theorem together with the statement about connected components, and
deduce that a strictly diagonally dominant matrix is invertible. Give an example showing that the
union of the discs cannot be replaced by their intersection with the corresponding column discs.

\item Prove the Courant--Fischer minimax theorem, deduce Cauchy's interlacing theorem, and derive
Weyl's inequality $\lambda_{i+j-1}(A+B)\leqslant\lambda_i(A)+\lambda_j(B)$ for Hermitian matrices.

\item (VJIMC 2003, cat.~I) Let $A$ and $B$ be complex Hermitian $2\times2$ matrices with eigenvalue
pairs $(\alpha_1,\alpha_2)$ and $(\beta_1,\beta_2)$. Determine all possible eigenvalue pairs
$(\gamma_1,\gamma_2)$ of $C=A+B$.

\item (VJIMC 2012, cat.~II) Let $M$ be the tridiagonal $10\times10$ matrix with $-1$ in the corner
$m_{11}$, with $3$ in the positions $m_{12}$ and $m_{21}$, with $2$ on the rest of the main
diagonal and $-1$ on the rest of the two neighbouring diagonals. Show that $M$ has exactly nine
positive real eigenvalues, counted with multiplicity.

\item (SEEMOUS 2012) Let $k$ be a positive integer and let $A$ be a real symmetric $n\times n$
matrix with $\bigl(\tr(A^k)\bigr)^{k+1}=\bigl(\tr(A^{k+1})\bigr)^k$.
\begin{enumerate}[label=\alph*),align=left]
\item Prove that $A^n=\tr(A)\,A^{n-1}$ when $k$ is even.
\item Does the same conclusion hold for odd $k$?
\end{enumerate}

\item (SEEMOUS 2018) Let $A,B\in M_{2018}(\mathbb R)$ satisfy $AB=BA$ and $A^{2018}=B^{2018}=E$.
Prove that if $\tr(AB)=2018$, then $\tr A=\tr B$.

\item (after IMC 2025 shortlist) Let $n$ be a positive integer. Find the least positive integer
$k$ with the following property: for every complex $n\times n$ matrix $A$ at least one of the $k$
matrices $E+A,\ E+A^2,\ \dots,\ E+A^k$ is invertible.

\item (IMC) Let $A$ and $B$ be invertible complex $n\times n$ matrices with $A+B=E$ and
$(A^2+B^2)(A^4+B^4)=A^5+B^5$. Find all possible values of $\det(AB)$.

\item (Romania national olympiad 2007) For $n\geqslant2$ let $H_n$ be the set of column vectors
$x\in\mathbb R^n$ with $\sum_{i=1}^n|x_i|=1$. Prove that there are only finitely many matrices
$A\in M_n(\mathbb R)$ for which the map $x\mapsto Ax$ satisfies $A(H_n)=H_n$.

\item Prove the Hoffman--Wielandt inequality
$\sum_i|\lambda_i(A)-\lambda_{\pi(i)}(B)|^2\leqslant\|A-B\|_F^2$ for normal $A,B$ and a suitable
permutation $\pi$, and deduce the standard bound
$\|\delta x\|/\|x\|\leqslant\varkappa(A)\,\|\delta b\|/\|b\|$ for perturbations of the system
$Ax=b$.
\end{enumerate}
\end{document}
